\chapter{HASIL DAN PEMBAHASAN}
\label{bab:hasil_pembahasan}

\section{Hasil Penelitian}
\label{sec:hasil}

Dummy konten hasil penelitian. Bagian ini menyajikan temuan dari eksperimen yang telah dilakukan.

\subsection{Hasil Eksperimen Model A}
\label{subsec:hasil_model_a}

Tabel \ref{tab:hasil_a} menunjukkan hasil eksperimen Model A.

\begin{table}[h]
    \centering
    \caption{Hasil Eksperimen Model A}
    \label{tab:hasil_a}
    \begin{tabularx}{\textwidth}{X c c c}
        \toprule
        \textbf{Dataset} & \textbf{Akurasi (\%)} & \textbf{F1-Score (\%)} & \textbf{Waktu (s)} \\
        \midrule
        Dataset 1 & 85.2 & 84.8 & 120 \\
        Dataset 2 & 87.5 & 87.1 & 150 \\
        Dataset 3 & 82.3 & 81.9 & 100 \\
        \bottomrule
    \end{tabularx}
\end{table}

\subsection{Hasil Eksperimen Model B}
\label{subsec:hasil_model_b}

Tabel \ref{tab:hasil_b} menunjukkan hasil eksperimen Model B.

\begin{table}[h]
    \centering
    \caption{Hasil Eksperimen Model B}
    \label{tab:hasil_b}
    \begin{tabularx}{\textwidth}{X c c c}
        \toprule
        \textbf{Dataset} & \textbf{Akurasi (\%)} & \textbf{F1-Score (\%)} & \textbf{Waktu (s)} \\
        \midrule
        Dataset 1 & 88.7 & 88.2 & 125 \\
        Dataset 2 & 90.1 & 89.6 & 155 \\
        Dataset 3 & 86.5 & 86.0 & 105 \\
        \bottomrule
    \end{tabularx}
\end{table}

\section{Analisis Komparatif}
\label{sec:analisis}

Dummy konten analisis komparatif. Bagian ini membandingkan hasil antara Model A dan Model B.

\subsection{Perbandingan Akurasi}
\label{subsec:perbandingan_akurasi}

Gambar \ref{fig:perbandingan} menunjukkan perbandingan akurasi antara kedua model.

\begin{figure}[h]
    \centering
    \fbox{\parbox[c][5cm][c]{12cm}{\centering Grafik Perbandingan Akurasi}}
    \caption{Perbandingan Akurasi Model A dan Model B}
    \label{fig:perbandingan}
\end{figure}

\subsection{Uji Signifikansi Statistik}
\label{subsec:uji_statistik}

Hasil uji statistik menggunakan paired t-test:

\begin{equation}
    t = \frac{\bar{d}}{s_d / \sqrt{n}}
    \label{eq:ttest}
\end{equation}

Dimana $\bar{d}$ adalah rata-rata perbedaan, $s_d$ adalah standar deviasi perbedaan, dan $n$ adalah jumlah sampel.

\section{Pembahasan}
\label{sec:pembahasan}

Dummy konten pembahasan. Bagian ini menginterpretasikan hasil penelitian dalam konteks teori dan penelitian sebelumnya.

\subsection{Interpretasi Hasil}
\label{subsec:interpretasi}

Interpretasi hasil penelitian dalam konteks hipotesis yang telah ditetapkan.

\subsection{Implikasi Teoretis}
\label{subsec:implikasi_teoretis}

Diskusi tentang kontribusi teoretis dari temuan penelitian.

\subsection{Implikasi Praktis}
\label{subsec:implikasi_praktis}

Diskusi tentang kontribusi praktis dari temuan penelitian.

\subsection{Keterbatasan}
\label{subsec:keterbatasan}

Diskusi tentang keterbatasan penelitian yang perlu diakui.